\section{See What Humans Want to See：把视觉能力约束到人的目标}

\subsection{从 ``replace humans'' 转向 ``augment humans''}

这一部分是整场讲座的价值核心。讲者明确反对把 AI 粗暴地理解为取代劳动的机器，相反，她反复强调的词是 \textbf{augment}。也就是说，系统的首要目标应当是\textbf{扩大人的能力边界、缓解稀缺人力、降低重复劳动负担、提升照护和决策覆盖范围}，而不是默认让机器接管一切。

这种立场并不天真。因为很多最重要的应用场景，比如医疗、养老、家庭照护，本来就存在巨大的人力缺口。AI 在这里不是和人争工作，而是在帮系统填补长期无法靠人力单独填平的空白。

\begin{importantbox}{Replace 与 Augment 的差别}
\begin{itemize}
    \item \textbf{Replace mindset}：把人看成成本，把机器看成更便宜替代品。
    \item \textbf{Augment mindset}：把人看成责任主体，把机器看成扩大覆盖面与一致性的工具。
\end{itemize}
Lecture 18 明显站在后者这一边。
\end{importantbox}

\subsection{医疗 ambient intelligence 是 human-centered AI 的代表场景}

前面提到的医疗监测，在这里被重新赋予了价值解释。为什么讲者如此看重 ambient intelligence？因为它天然符合 human-centered design 的四个条件：

\begin{enumerate}
    \item \textbf{需求真实}：病房、ICU、家庭照护都有持续监测需求。
    \item \textbf{人力稀缺}：不可能全天候依赖人工高质量覆盖。
    \item \textbf{风险可定义}：跌倒、异常活动、物品遗漏、手卫生等都有明确目标。
    \item \textbf{责任边界清楚}：AI 给出信号，人类负责最终理解与处置。
\end{enumerate}

这也是 lecture18 的重要判断：\textbf{最值得优先做的 AI，不一定是最炫的，而是最能在真实世界里减轻人的痛点、同时保留人类主导权的。}

\subsection{视觉的终点不是描述世界，而是进入行动闭环}

如果说前两个部分解决的是``看见什么''与``怎么看得更好''，那么这一部分开始处理``看见之后如何行动''。讲者把话题推进到 embodied AI 与机器人，这意味着视觉系统不再停留在认知层，而是要参与决策、控制和任务完成。

机器人场景里的难点显著升级：

\begin{itemize}
    \item 环境开放且变化剧烈，不再是静态 benchmark。
    \item 任务是组合式的，涉及多步骤依赖。
    \item 感知错误会直接传递到动作错误。
    \item 评估标准不再只是 accuracy，而是任务完成率、鲁棒性与安全性。
\end{itemize}

讲者给出的框架很有代表性：大模型负责把人类意图转成子任务，视觉模型理解当前场景，规划模块把这些信息变成可执行的 action sequence，最终驱动机器人操作。这里最重要的不是某一个模块，而是 perception--language--planning--action 的闭环。

\begin{figure}[H]
\centering
\begin{tikzpicture}[every node/.style={font=\small, align=center}]
\node[draw, rounded corners, fill=green!8, minimum width=2.2cm, minimum height=0.95cm] (goal) at (0,0) {人类目标\\task request};
\node[draw, rounded corners, fill=green!12, minimum width=2.4cm, minimum height=0.95cm] (vlm) at (3.2,0) {视觉理解\\objects, state};
\node[draw, rounded corners, fill=green!16, minimum width=2.6cm, minimum height=0.95cm] (llm) at (6.8,0) {任务分解\\subgoals / plan};
\node[draw, rounded corners, fill=green!20, minimum width=2.6cm, minimum height=0.95cm] (motion) at (10.4,0) {动作规划\\motion / control};
\node[draw, rounded corners, fill=green!10, minimum width=2.6cm, minimum height=0.95cm] (robot) at (13.8,0) {机器人执行\\feedback};
\draw[-{Latex[length=3mm]}, thick] (goal) -- (vlm);
\draw[-{Latex[length=3mm]}, thick] (vlm) -- (llm);
\draw[-{Latex[length=3mm]}, thick] (llm) -- (motion);
\draw[-{Latex[length=3mm]}, thick] (motion) -- (robot);
\end{tikzpicture}
\caption{具身 AI 闭环：目标、理解、分解、规划、执行}
\end{figure}

\begin{knowledgebox}{为什么机器人会逼迫视觉系统更真实}
因为机器人不能只在 representation space 里``看起来懂了''。一旦要落到执行，错误就会通过抓取失败、路径碰撞、任务中断等形式被立即暴露。机器人因此天然要求视觉表示更稳定、更结构化、更可操作。
\end{knowledgebox}

\subsection{Open-World Robotics：LLM 和 VLM 是怎么接到机器人上的}

这部分字幕给出的不是抽象愿景，而是一条相当具体的技术路线。讲者想解决的问题是：\textbf{机器人能不能像 LLM 一样接受开放指令，而不是只能在预定义闭世界里执行少数固定任务？}\footnote{根据字幕 00:49:21--00:53:18，讲者介绍 open instruction robotics：LLM 生成代码或子任务，VLM 识别环境中的 drawer / handle / vase，再把结果更新到 motion planning heatmap 中。}

她给出的例子是 ``open top drawer, but avoid the vase''。要完成这件事，系统至少要做四层映射：

\begin{enumerate}
    \item 把自然语言指令解析成可执行子目标。
    \item 在当前场景里找到 drawer、handle、vase 等目标物。
    \item 把 ``靠近抽屉把手'' 和 ``远离花瓶'' 共同转化成 motion planning map。
    \item 将空间热图进一步落到 gripper 旋转、速度和路径控制上。
\end{enumerate}

\begin{figure}[H]
\centering
\begin{tikzpicture}[every node/.style={font=\small, align=center}]
\node[draw, rounded corners, fill=orange!10, minimum width=2.2cm, minimum height=0.9cm] (text) at (0,0) {开放指令\\open top drawer};
\node[draw, rounded corners, fill=orange!16, minimum width=2.5cm, minimum height=0.9cm] (llm) at (3.4,0) {LLM\\子任务/代码生成};
\node[draw, rounded corners, fill=orange!22, minimum width=2.5cm, minimum height=0.9cm] (vlm) at (7.0,0) {VLM\\drawer / handle / vase};
\node[draw, rounded corners, fill=orange!28, minimum width=2.8cm, minimum height=0.9cm] (map) at (10.9,0) {motion heatmap\\where to go / avoid};
\node[draw, rounded corners, fill=orange!14, minimum width=2.4cm, minimum height=0.9cm] (ctrl) at (14.5,0) {control\\trajectory / gripper};
\draw[-{Latex[length=3mm]}, thick] (text) -- (llm);
\draw[-{Latex[length=3mm]}, thick] (llm) -- (vlm);
\draw[-{Latex[length=3mm]}, thick] (vlm) -- (map);
\draw[-{Latex[length=3mm]}, thick] (map) -- (ctrl);
\end{tikzpicture}
\caption{open-world robotics 管线：语言 grounding 后进入动作规划}
\end{figure}

这条路线真正解决的，不是某个特定 drawer 的开合，而是机器人``泛化到没见过的环境''的能力。以前很多机器人系统之所以显得脆弱，是因为它们在训练时就已经知道世界会长成什么样；而这套 LLM + VLM + planner 的组合，至少提供了一种离开 closed world 的办法。

\begin{warningbox}{为什么这还远远没有被解决}
因为真实家庭环境里，目标会遮挡、光照会变、物体会移动、指令会含糊、动作后果会连锁变化。LLM 和 VLM 能帮助开放指令与开放感知，但它们还没有把机器人的可靠性问题彻底解决。
\end{warningbox}

\subsection{Benchmark 也必须从 ``漂亮'' 走向 ``生态有效''}

讲者提到 ecological robotic learning environment 与 BEHAVIOR benchmark，重点不在于介绍一个数据集名称，而在于提出 benchmark 设计哲学的变化：\textbf{如果评测环境离真实家庭、真实物体、真实多步骤任务太远，那么机器人系统很可能只是在刷分，而没有变得真正有用。}\footnote{根据字幕 00:54:40--00:55:10，讲者介绍 ecological robotic learning 与 BEHAVIOR benchmark，强调 everyday household activity 与 virtual interactive ecological environments。}

这和课程前面讲的 ImageNet 时代形成了呼应。早期 benchmark 的价值是把识别问题规模化、标准化；而到了机器人和 agent 阶段，benchmark 的价值变成了让模型在更接近真实世界的组合任务里暴露真实短板。

\begin{warningbox}{Benchmark 的陷阱}
如果 benchmark 只奖励单步成功率、固定视角和干净场景，那么模型会学会迎合评测，而不是学会服务真实世界。Lecture 18 要求我们重新思考：什么样的 benchmark 才配得上 human-centered AI。
\end{warningbox}

\subsection{为什么 BEHAVIOR 这类 benchmark 会让人 ``有点沮丧''}

讲者后面给出的数字其实很刺耳：在不给 privileged information 的前提下，把今天的算法直接丢进这类生态任务环境，很多任务的表现几乎就是 \textbf{0}。\footnote{根据字幕 01:00:30--01:01:29，讲者明确说 today's algorithms still cannot do behavior tasks；在没有 privileged information 的 top row 条件下，三个 behavior task 的表现接近 zero。}

这不是坏消息，反而是 benchmark 真正有价值的证据。因为如果一个 benchmark 让大家轻松刷出高分，那往往只能说明任务过于理想化；而当系统一旦失去完美记忆、完美状态、magic motion 之类的强假设，性能就迅速掉到接近零，说明 benchmark 终于开始逼近真实世界的复杂性。

\begin{importantbox}{生态 benchmark 暴露出的真实缺口}
\begin{itemize}
    \item 当前算法离长程规划仍然很远。
    \item 当前算法离真实物理交互仍然很远。
    \item 当前算法离开放环境中的稳健泛化仍然很远。
    \item 当前算法离``没有特权信息也能做事''仍然很远。
\end{itemize}
\end{importantbox}

讲者还提到，BEHAVIOR 不只是一个虚拟环境，而是在 50 个真实场景扫描、上万物体资产、变形体/透明体/热学等属性建模、以及与 Omniverse 协同构建的物理与感知仿真基础上发展出来的。这意味着它追求的不只是 render 得好看，而是让任务结构、物理属性和感知输入都尽量接近现实。

\subsection{从数字孪生到脑机接口：lecture18 给出的未来方向}

在收尾前，讲者还快速扫过了几条未来线索：real-to-sim transfer、digital twin、用同一生态环境研究视觉障碍患者，以及利用 EEG 控制机器人完成烹饪等动作。\footnote{根据字幕 01:01:44--01:04:29，讲者提到数字孪生、视觉障碍研究场景，以及通过 EEG 信号控制机器人完成整套操作。}

这些例子共同指向一件事：视觉 AI 的未来不是只做 ``更强分类器''，而是和机器人、临床、认知科学、辅助技术深度耦合。尤其是 EEG 控制机器人的例子，虽然还非常早期，但它很直观地展示了 lecture18 想强调的方向：\textbf{AI 与机器人最动人的用途之一，是帮助那些因为身体条件受限而无法直接行动的人恢复行动能力。}

\subsection{决定机器人该做什么的人，不应该只是研究者}

这一段是 lecture18 最``人本''、也最不像传统技术课的一段。讲者提出一个很直接的问题：\textbf{到底是谁来决定机器人应该学哪些任务？} 如果只靠实验室研究者拍脑袋，得到的任务列表很可能只是``技术上好做''，或者``研究者自己觉得有趣''，而不一定是普通人真正愿意交给机器的事情。

因此，讲者介绍了一个 human-centered survey：研究者汇总了来自美国和欧洲劳工调查中的大量日常活动任务，再邀请约 1,400 名受访者表达他们愿意让机器人帮忙完成哪些工作、又不愿意把哪些任务交出去。\footnote{根据字幕 00:55:10--00:57:18，讲者介绍从政府劳动调查中整理日常任务，并在线招募约 1,400 名参与者评估机器人任务偏好。}

这个结果非常有信息量。人们一般欢迎机器人做清洁、铲雪、重复劳动、脏活累活；但对情感性强、社交意义强、象征意义强的任务则明显抗拒，比如打开圣诞礼物、购买婚戒、某些婴儿照护细节等。也就是说，\textbf{任务是否自动化，不只取决于技术可行性，还取决于情感、社会与文化含义。}

\begin{figure}[H]
\centering
\begin{tikzpicture}[font=\small]
\draw[->] (0,0) -- (12,0);
\node[anchor=east] at (-0.2,2.8) {偏好高};
\node[anchor=east] at (-0.2,2.0) {偏好中};
\node[anchor=east] at (-0.2,1.2) {偏好低};
\fill[blue!45] (0.6,2.5) rectangle (4.6,2.9);
\fill[blue!35] (0.6,1.7) rectangle (3.8,2.1);
\fill[blue!25] (0.6,0.9) rectangle (2.4,1.3);
\fill[orange!55] (6.2,2.5) rectangle (7.8,2.9);
\fill[orange!45] (6.2,1.7) rectangle (8.7,2.1);
\fill[orange!35] (6.2,0.9) rectangle (10.8,1.3);
\node[anchor=west] at (4.8,2.7) {厨房/地面清洁};
\node[anchor=west] at (4.0,1.9) {铲雪、洗碗、折衣物};
\node[anchor=west] at (2.6,1.1) {做早餐等混合任务};
\node[anchor=west] at (8.0,2.7) {娱乐性或象征性任务};
\node[anchor=west] at (8.9,1.9) {打开礼物、买婚戒};
\node[anchor=west] at (11.0,1.1) {情感/社交意义强的照护行为};
\end{tikzpicture}
\caption{机器人任务偏好：脏活累活偏好高，情感与象征性任务偏好低\protect\footnotemark}
\end{figure}
\footnotetext{根据字幕 00:55:10--00:58:01 整理。讲者给出的实例包括 clean kitchen floor、shoveling snow、folding laundry、opening Christmas gift、buying wedding rings 等。}

\begin{importantbox}{这项调查的价值不在于给出一个固定榜单}
它真正重要的地方在于提出了一种原则：\textbf{机器人任务空间不该只由技术供给决定，也要由人的偏好来定义。} Human-centered AI 不是在系统做完之后征求意见，而是在任务定义阶段就把人的选择权纳入进来。
\end{importantbox}

\subsection{这一部分真正想传达什么}

``See what humans want to see'' 的意思不是``让 AI 更讨喜''，而是把系统从根上绑定到人的目的上。它要求研究者同时回答三个问题：

\begin{enumerate}
    \item 这个系统究竟在帮谁？
    \item 它是在替代人的判断，还是在增强人的能力？
    \item 它优化的任务目标，是技术人员自定的，还是经过人的真实偏好校验的？
\end{enumerate}

\begin{knowledgebox}{Lecture 18 的核心价值判断}
真正高级的 AI 不是能力无限扩张，而是在能力越来越强时，仍然知道自己该停在哪、该向谁负责、该为哪些真实需求服务。
\end{knowledgebox}

